\documentclass[11pt,a4paper]{article}

\usepackage[margin=1in]{geometry}
\usepackage[T1]{fontenc}
\usepackage[hidelinks,colorlinks=true,linkcolor=black,urlcolor=blue,citecolor=black]{hyperref}
\usepackage{xcolor}
\usepackage{enumitem}
\usepackage{titlesec}
\usepackage{parskip}

\definecolor{brandbrown}{HTML}{A9704F}
\titleformat{\section}{\normalfont\Large\bfseries\color{brandbrown}}{\thesection}{1em}{}
\titleformat{\subsection}{\normalfont\large\bfseries}{\thesubsection}{1em}{}

\title{\vspace{-1cm}\textbf{Computer Science Portfolio}\\[4pt]\large Project Report}
\author{Maide Oskay}
\date{September 2026}

\begin{document}
\maketitle
\vspace{-0.5cm}

\section*{Links}
\begin{itemize}[leftmargin=1.2cm]
  \item \textbf{Portfolio website:} \url{https://maidesource.github.io}
  \item \textbf{GitHub repository:} \url{https://github.com/maidesource/maidesource.github.io}
\end{itemize}

\section{Introduction}
This report accompanies the computer science portfolio I built for a university course
assignment, which simulates preparing a portfolio for a working student position application.
The portfolio consists of a personal website hosted with GitHub Pages
and its corresponding public GitHub repository, which together present my education, technical
skills, and ongoing projects, alongside a LaTeX-typeset curriculum vitae. This document explains
how the portfolio is structured, the reasoning behind its design, the tools used to build it, and
a reflection on its current strengths, limitations, and planned improvements.

\section{Repository and Website Structure}
The repository is named \texttt{maidesource.github.io}, which GitHub automatically recognises as
a user site and serves directly from the \texttt{main} branch at
\url{https://maidesource.github.io}. The repository root is kept flat and predictable:

\begin{itemize}
  \item \texttt{index.html}: the single-page website, containing markup for the header, About,
        Education, Projects, and Skills sections, and the footer.
  \item \texttt{styles.css}: all visual styling, including the color palette, typography,
        spacing, and the card-based layout used for project entries.
  \item \texttt{script.js}: a small script that enables smooth scrolling when the navigation
        links are clicked.
  \item \texttt{assets/}: the profile photograph used in the page header.
  \item \texttt{docs/}: the compiled CV (\texttt{cv.pdf}) together with its LaTeX source
        (\texttt{cv.tex}), so the CV can be regenerated or edited without leaving the repository.
  \item \texttt{projects/}: documentation of the technical exercise showcased in the
        portfolio. At this stage, the website itself is treated as the featured project, and this
        folder explains why and how it was built (see \texttt{projects/README.md}).
  \item \texttt{README.md}: a top-level description of the repository's purpose and layout for
        anyone visiting the project on GitHub.
\end{itemize}

On the website itself, the navigation bar links to in-page sections (About, Projects, Skills)
using anchor links, and to the CV, which opens the PDF stored in \texttt{docs/}. This keeps the
whole experience of browsing the site and reading the CV inside a single, coherent repository,
rather than splitting content across external services.

\section{Design Decisions}
Several deliberate choices shaped the final result:

\begin{itemize}
  \item \textbf{Single-page layout.} A hiring manager should be able to see education, skills,
        and projects within a few seconds of scrolling, without navigating through multiple
        pages. A single page with anchored sections keeps the site quick to scan while still
        feeling organised.
  \item \textbf{Consistent visual identity.} The color palette (warm terracotta and beige tones)
        was carried over from my personal CV design so that the website and CV are immediately
        recognisable as belonging to the same person and brand.
  \item \textbf{Card-based project entries.} Each project is presented in its own visually
        distinct card with a title, date range, and bullet points. This mirrors how the same
        content is organised in the CV, so a reader who has seen one recognises the same
        information in the other.
  \item \textbf{Minimal dependencies.} The site intentionally avoids CSS or JavaScript
        frameworks. For a page of this size and complexity, a framework would add
        download weight and hide the underlying HTML/CSS/JS skills that the portfolio is meant to
        demonstrate in the first place.
  \item \textbf{Version-controlled CV.} Rather than treating the CV as an external, opaque PDF,
        its LaTeX source lives in the same repository. This makes the CV auditable, diffable, and
        easy to update the same way the rest of the site is updated: edit, commit, push.
\end{itemize}

\section{Tools and Technologies}
\begin{itemize}
  \item \textbf{HTML5 and CSS3} for the page structure and styling, using semantic elements
        (\texttt{header}, \texttt{main}, \texttt{section}, \texttt{footer}) and Flexbox for
        layout.
  \item \textbf{Vanilla JavaScript (ES6)} for the smooth-scroll navigation behaviour, with no
        external libraries.
  \item \textbf{Git and GitHub} for version control, and \textbf{GitHub Pages} for hosting and
        continuous deployment: every push to \texttt{main} automatically republishes the live
        site.
  \item \textbf{LaTeX} (compiled with \texttt{pdflatex} on TeX Live) for typesetting the CV,
        allowing precise control over layout and easy long-term maintenance as a plain-text
        source file.
  \item A Linux terminal (Fedora) with \texttt{nano} was used throughout for editing files and
        running Git commands, reinforcing familiarity with a command-line development workflow.
\end{itemize}

\section{Reflection}

\subsection{Strengths}
\begin{itemize}
  \item The site is lightweight and loads quickly, since it has no external framework
        dependencies.
  \item Content, styling, and behaviour are cleanly separated into
        \texttt{index.html}/\texttt{styles.css}/\texttt{script.js}, which keeps the codebase easy
        to read and extend.
  \item The CV is reproducible and version-controlled rather than being a static, unexplained
        binary file.
  \item The visual identity is consistent across the website and CV, giving a professional and
        cohesive impression.
  \item The deployment pipeline is simple and reliable: GitHub Pages removes the need for any
        separate hosting setup.
\end{itemize}

\subsection{Weaknesses}
\begin{itemize}
  \item The layout is not yet fully responsive; it has not been adapted with media queries for
        small screens, so the experience on mobile devices is not optimised.
  \item The \texttt{projects/} folder currently documents the portfolio site itself rather than a
        separate body of coding exercises, since I do not yet have additional finished technical
        projects to include.
  \item There is no automated testing, linting, or build process; the site is edited and deployed
        manually.
  \item No accessibility audit (colour contrast, screen-reader behaviour) has been performed yet.
\end{itemize}

\subsection{Planned Improvements}
\begin{itemize}
  \item Add responsive, mobile-first CSS using media queries so the site adapts cleanly to
        smaller screens.
  \item Populate \texttt{projects/} with additional, independent technical exercises (small
        scripts, algorithms, or coursework projects) as they are completed.
  \item Introduce basic automated checks (e.g. an HTML/CSS validator or a Lighthouse
        accessibility and performance audit) as part of the development workflow.
  \item Improve SEO and shareability with meta tags such as an Open Graph description and page
        summary.
  \item Consider a short "notes" or "learning log" section documenting technical topics studied
        during my degree, to keep the portfolio active between larger projects.
\end{itemize}

\section{Conclusion}
The portfolio in its current form fulfils its core purpose: it presents my background, skills,
and ongoing projects clearly, is fully version-controlled, and is publicly accessible through
GitHub Pages. Its main limitations, namely a thin project section and a non-responsive layout,
are well understood and have a concrete improvement plan, which I intend to carry out as I
complete further coursework and personal projects.

\end{document}
